\subsection{CFT's From Calabi-Yau Four-folds --- arXiv:hep-th/9906070}

\textbf{Authors:} Sergei Gukov, Cumrun Vafa, Edward Witten

\paragraph{主な主張}
フラックスとブレーンを含むCalabi--Yau四次元多様体の孤立特異点近傍を解析し、真空・ソリトン構造から有効超ポテンシャルや低次元の非自明な超共形理論が現れることを示した。 \cite{Gukov:1999ya}

\paragraph{新規性}
フラックスと正則体積形式から超ポテンシャルを記述し、ブレーンが作るdomain wallの張力と超ポテンシャル差を対応付けた。

\paragraph{位置付け}
Gukov--Vafa--Witten（GVW）フラックス超ポテンシャルの基礎文献であり、後のフラックスによるモジュライ安定化を支える。
